\documentclass[10pt,a4paper]{amsart}
\usepackage{amsmath,amssymb,amsthm,graphicx}
\usepackage{verbatim}
\newenvironment{reference}{\comment}{\endcomment}
%\newenvironment{reference}{}{}
\newenvironment{slogan}{\comment}{\endcomment}
\newenvironment{history}{\comment}{\endcomment}

% For commutative diagrams we use Xy-pic
\usepackage[all]{xy}

% We use 2cell for 2-commutative diagrams.
\xyoption{2cell}
% Theorem environments.
%
\theoremstyle{plain}
\newtheorem{theorem}[subsection]{Theorem}
\newtheorem{proposition}[subsection]{Proposition}
\newtheorem{lemma}[subsection]{Lemma}

\theoremstyle{definition}
\newtheorem{definition}[subsection]{Definition}
\newtheorem{example}[subsection]{Example}
\newtheorem{exercise}[subsection]{Exercise}
\newtheorem{situation}[subsection]{Situation}

\theoremstyle{remark}
\newtheorem{remark}[subsection]{Remark}
\newtheorem{remarks}[subsection]{Remarks}

\numberwithin{equation}{subsection}

% Macros
%
\def\lim{\mathop{\mathrm{lim}}\nolimits}
\def\colim{\mathop{\mathrm{colim}}\nolimits}
\def\Spec{\mathop{\mathrm{Spec}}}
\def\Hom{\mathop{\mathrm{Hom}}\nolimits}
\def\Ext{\mathop{\mathrm{Ext}}\nolimits}
\def\SheafHom{\mathop{\mathcal{H}\!\mathit{om}}\nolimits}
\def\SheafExt{\mathop{\mathcal{E}\!\mathit{xt}}\nolimits}
\def\Sch{\mathit{Sch}}
\def\Mor{\mathop{\mathrm{Mor}}\nolimits}
\def\Ob{\mathop{\mathrm{Ob}}\nolimits}
\def\Sh{\mathop{\mathit{Sh}}\nolimits}
\def\NL{\mathop{N\!L}\nolimits}
\def\CH{\mathop{\mathrm{CH}}\nolimits}
\def\proetale{{pro\text{-}\acute{e}tale}}
\def\etale{{\acute{e}tale}}
\def\QCoh{\mathit{QCoh}}
\def\Ker{\mathop{\mathrm{Ker}}}
\def\Im{\mathop{\mathrm{Im}}}
\def\Coker{\mathop{\mathrm{Coker}}}
\def\Coim{\mathop{\mathrm{Coim}}}

% Boxtimes
%
\DeclareMathSymbol{\boxtimes}{\mathbin}{AMSa}{"02}

%
% Macros for moduli stacks/spaces
%
\def\QCohstack{\mathcal{QC}\!\mathit{oh}}
\def\Cohstack{\mathcal{C}\!\mathit{oh}}
\def\Spacesstack{\mathcal{S}\!\mathit{paces}}
\def\Quotfunctor{\mathrm{Quot}}
\def\Hilbfunctor{\mathrm{Hilb}}
\def\Curvesstack{\mathcal{C}\!\mathit{urves}}
\def\Polarizedstack{\mathcal{P}\!\mathit{olarized}}
\def\Complexesstack{\mathcal{C}\!\mathit{omplexes}}
% \Pic is the operator that assigns to X its picard group, usage \Pic(X)
% \Picardstack_{X/B} denotes the Picard stack of X over B
% \Picardfunctor_{X/B} denotes the Picard functor of X over B
\def\Pic{\mathop{\mathrm{Pic}}\nolimits}
\def\Picardstack{\mathcal{P}\!\mathit{ic}}
\def\Picardfunctor{\mathrm{Pic}}
\def\Deformationcategory{\mathcal{D}\!\mathit{ef}}
\usepackage{cleveref}
\crefname{originalreference}{original reference}{original references}
\Crefname{originalreference}{Original reference}{Original references}
\makeatletter
\newlabel{situation-shriek}{{0F42}{source}{}{sourcepointer-1}{}}
\newlabel{situation-shriek@cref}{{[originalreference][0][]0F42}{[1][0][]source}{}{}{}}
\makeatother
\crefname{originallemma}{original lemma}{original lemmas}
\Crefname{originallemma}{Original lemma}{Original lemmas}
\makeatletter
\newlabel{lemma-map-pullback-to-shriek-well-defined}{{0B6T}{source}{}{sourcepointer-2}{}}
\newlabel{lemma-map-pullback-to-shriek-well-defined@cref}{{[originallemma][0][]0B6T}{[1][0][]source}{}{}{}}
\makeatother
\crefname{originallemma}{original lemma}{original lemmas}
\Crefname{originallemma}{Original lemma}{Original lemmas}
\makeatletter
\newlabel{more-morphisms-lemma-flat-finite-presentation-perfect}{{068A}{source}{}{sourcepointer-3}{}}
\newlabel{more-morphisms-lemma-flat-finite-presentation-perfect@cref}{{[originallemma][0][]068A}{[1][0][]source}{}{}{}}
\makeatother
\crefname{originalreference}{original reference}{original references}
\Crefname{originalreference}{Original reference}{Original references}
\makeatletter
\newlabel{remark-local-calculation-shriek}{{0BV2}{source}{}{sourcepointer-4}{}}
\newlabel{remark-local-calculation-shriek@cref}{{[originalreference][0][]0BV2}{[1][0][]source}{}{}{}}
\makeatother
\crefname{originallemma}{original lemma}{original lemmas}
\Crefname{originallemma}{Original lemma}{Original lemmas}
\makeatletter
\newlabel{dualizing-lemma-relative-dualizing-algebraic}{{0BZW}{source}{}{sourcepointer-5}{}}
\newlabel{dualizing-lemma-relative-dualizing-algebraic@cref}{{[originallemma][0][]0BZW}{[1][0][]source}{}{}{}}
\makeatother
\crefname{originallemma}{original lemma}{original lemmas}
\Crefname{originallemma}{Original lemma}{Original lemmas}
\makeatletter
\newlabel{perfect-lemma-tor-dimension-rel-affine}{{0DHY}{source}{}{sourcepointer-6}{}}
\newlabel{perfect-lemma-tor-dimension-rel-affine@cref}{{[originallemma][0][]0DHY}{[1][0][]source}{}{}{}}
\makeatother
\crefname{originallemma}{original lemma}{original lemmas}
\Crefname{originallemma}{Original lemma}{Original lemmas}
\makeatletter
\newlabel{perfect-lemma-pseudo-coherent}{{08E5}{source}{}{sourcepointer-7}{}}
\newlabel{perfect-lemma-pseudo-coherent@cref}{{[originallemma][0][]08E5}{[1][0][]source}{}{}{}}
\makeatother
\crefname{originallemma}{original lemma}{original lemmas}
\Crefname{originallemma}{Original lemma}{Original lemmas}
\makeatletter
\newlabel{perfect-lemma-identify-pseudo-coherent-noetherian}{{08E8}{source}{}{sourcepointer-8}{}}
\newlabel{perfect-lemma-identify-pseudo-coherent-noetherian@cref}{{[originallemma][0][]08E8}{[1][0][]source}{}{}{}}
\makeatother
\crefname{originallemma}{original lemma}{original lemmas}
\Crefname{originallemma}{Original lemma}{Original lemmas}
\makeatletter
\newlabel{perfect-lemma-quasi-coherence-internal-hom}{{0A6H}{source}{}{sourcepointer-9}{}}
\newlabel{perfect-lemma-quasi-coherence-internal-hom@cref}{{[originallemma][0][]0A6H}{[1][0][]source}{}{}{}}
\makeatother
\crefname{originallemma}{original lemma}{original lemmas}
\Crefname{originallemma}{Original lemma}{Original lemmas}
\makeatletter
\newlabel{lemma-shriek-coherent}{{0AU1}{source}{}{sourcepointer-10}{}}
\newlabel{lemma-shriek-coherent@cref}{{[originallemma][0][]0AU1}{[1][0][]source}{}{}{}}
\makeatother
\crefname{originallemma}{original lemma}{original lemmas}
\Crefname{originallemma}{Original lemma}{Original lemmas}
\makeatletter
\newlabel{lemma-shriek-open-immersion}{{0AU0}{source}{}{sourcepointer-11}{}}
\newlabel{lemma-shriek-open-immersion@cref}{{[originallemma][0][]0AU0}{[1][0][]source}{}{}{}}
\makeatother
\crefname{originallemma}{original lemma}{original lemmas}
\Crefname{originallemma}{Original lemma}{Original lemmas}
\makeatletter
\newlabel{lemma-compare-with-pullback-flat-proper-noetherian}{{0A9U}{source}{}{sourcepointer-12}{}}
\newlabel{lemma-compare-with-pullback-flat-proper-noetherian@cref}{{[originallemma][0][]0A9U}{[1][0][]source}{}{}{}}
\makeatother
\crefname{originallemma}{original lemma}{original lemmas}
\Crefname{originallemma}{Original lemma}{Original lemmas}
\makeatletter
\newlabel{morphisms-lemma-quasi-affine-finite-type-over-S}{{04II}{source}{}{sourcepointer-13}{}}
\newlabel{morphisms-lemma-quasi-affine-finite-type-over-S@cref}{{[originallemma][0][]04II}{[1][0][]source}{}{}{}}
\makeatother
\crefname{originallemma}{original lemma}{original lemmas}
\Crefname{originallemma}{Original lemma}{Original lemmas}
\makeatletter
\newlabel{more-morphisms-lemma-perfect-permanence}{{068D}{source}{}{sourcepointer-14}{}}
\newlabel{more-morphisms-lemma-perfect-permanence@cref}{{[originallemma][0][]068D}{[1][0][]source}{}{}{}}
\makeatother
\crefname{originallemma}{original lemma}{original lemmas}
\Crefname{originallemma}{Original lemma}{Original lemmas}
\makeatletter
\newlabel{lemma-shriek-closed-immersion}{{0AA2}{source}{}{sourcepointer-15}{}}
\newlabel{lemma-shriek-closed-immersion@cref}{{[originallemma][0][]0AA2}{[1][0][]source}{}{}{}}
\makeatother
\crefname{originallemma}{original lemma}{original lemmas}
\Crefname{originallemma}{Original lemma}{Original lemmas}
\makeatletter
\newlabel{lemma-shriek-affine-line}{{0AA1}{source}{}{sourcepointer-16}{}}
\newlabel{lemma-shriek-affine-line@cref}{{[originallemma][0][]0AA1}{[1][0][]source}{}{}{}}
\makeatother
\title{For a perfect morphism the shriek of the structure sheaf is bounded coherent, of finite tor dimension, self dual, preserves bounded coherent objects and makes the comparison map an isomorphism}
\author{The Stacks Project Authors}
\date{stacks-project:duality.tex:local-snapshot-20261012; GNU FDL 1.2 or later}
\begin{document}
\maketitle
\noindent\textit{Editorial note.} All delivered mathematical text is the author's own, copied from the pinned Stacks Project chapter file duality.tex; no statement or proof step was rewritten, reordered or removed.\par\medskip
\noindent\textit{Editorial note.} Preamble provenance: the pinned chapter file opens with an input of the shared preamble and carries no preamble text of its own, so the project's single Stacks-wide preamble staged at sources/stacks-duality/preamble.tex is excerpted verbatim under the recipe fixed by LANE\_BRIEF 22.1, as three spans, lines 8-13 (verbatim package and the reference/slogan/history environments), lines 14-18 (Xy-pic setup) and lines 154-217 (the native theorem declarations and mathematical macros). Nothing was added, removed or redefined; the wrapper class and the amsmath/amssymb/amsthm/graphicx base packages are supplied by the assembler.\par\medskip
\noindent\textit{Editorial note.} The counted statement and proof contain no cross reference, no citation and no included file, so the delivered body is self contained.\par\medskip
\noindent\textit{Editorial note.} Canonical location of the counted proposition: Stacks tag 0B6U (sources/stacks-duality/pinned-original-source.tex line 4798).\par\medskip
\noindent\textit{Allowed standard prerequisites (editorial).}\begin{itemize}
\item The upper shriek of a closed immersion, of an open immersion and of the affine line projection. Same chapter, Lemmas labelled lemma-shriek-closed-immersion, lemma-shriek-open-immersion and lemma-shriek-affine-line, cited at proof lines 4844-4846, 4884-4886 and 4900-4902. Compute f\textasciicircum{}! in each of the model cases used in the reduction.
\item The comparison map mu and its compatibility with composition. Same chapter, Lemma labelled lemma-map-pullback-to-shriek-well-defined, invoked in the statement and cited at proof lines 4862-4868. Supplies the map whose bijectivity is assertion (e) and the diagram closing (e) under composition.
\item Cohomological bounds for the shriek of coherent complexes and the perfect proper case. Same chapter, Lemma labelled lemma-shriek-coherent, cited at proof lines 4840-4842, and Lemma labelled lemma-compare-with-pullback-flat-proper-noetherian, cited at proof lines 4852-4854. Give the first preliminary remark and settle (e) for a perfect proper morphism.
\item The local algebraic computation of the relative dualizing complex. Same chapter, Remark labelled remark-local-calculation-shriek and Dualizing Complexes, Lemma labelled dualizing-lemma-relative-dualizing-algebraic, cited at proof lines 4828-4832. Proves (a), (b) and (c) in the affine case.
\item Finiteness and perfection permanence for the factorization. Morphisms, Lemma labelled morphisms-lemma-quasi-affine-finite-type-over-S, cited at proof lines 4874-4876, and More on Morphisms, Lemmas labelled more-morphisms-lemma-flat-finite-presentation-perfect and more-morphisms-lemma-perfect-permanence, cited at proof lines 4819-4821 and 4878-4880. Produce the factorization into a perfect closed immersion and a projection.
\end{itemize}\medskip
\ifdefined\hypertarget\hypertarget{sourcepointer-1}{}\fi
\noindent\textit{Source reference.} Original reference 0F42: sources/stacks-duality/pinned-original-source.tex, source lines 3953--3953. Same-chapter pointer into the author environment situation at source line 3953; invoked at proof line None of the counted proof. It is retained only as a navigation pointer; every load-bearing step is inside the delivered span.\par\medskip
\ifdefined\hypertarget\hypertarget{sourcepointer-2}{}\fi
\noindent\textit{Source reference.} Original lemma 0B6T: sources/stacks-duality/pinned-original-source.tex, source lines 4427--4427. Same-chapter pointer into the author environment lemma at source line 4427; invoked at proof line 4855 of the counted proof. It is retained only as a navigation pointer; every load-bearing step is inside the delivered span.\par\medskip
\ifdefined\hypertarget\hypertarget{sourcepointer-3}{}\fi
\noindent\textit{Source reference.} Original lemma 068A: sources/stacks-chow/tags.txt, source lines 7999--7999. Cross-chapter standard prerequisite located in the project tag list staged in the local snapshot; invoked at proof line 4823 of the counted proof. It is retained only as a navigation pointer; every load-bearing step is inside the delivered span.\par\medskip
\ifdefined\hypertarget\hypertarget{sourcepointer-4}{}\fi
\noindent\textit{Source reference.} Original reference 0BV2: sources/stacks-duality/pinned-original-source.tex, source lines 4636--4636. Same-chapter pointer into the author environment remark at source line 4636; invoked at proof line 4826 of the counted proof. It is retained only as a navigation pointer; every load-bearing step is inside the delivered span.\par\medskip
\ifdefined\hypertarget\hypertarget{sourcepointer-5}{}\fi
\noindent\textit{Source reference.} Original lemma 0BZW: sources/stacks-chow/tags.txt, source lines 15066--15066. Cross-chapter standard prerequisite located in the project tag list staged in the local snapshot; invoked at proof line 4828 of the counted proof. It is retained only as a navigation pointer; every load-bearing step is inside the delivered span.\par\medskip
\ifdefined\hypertarget\hypertarget{sourcepointer-6}{}\fi
\noindent\textit{Source reference.} Original lemma 0DHY: sources/stacks-chow/tags.txt, source lines 16923--16923. Cross-chapter standard prerequisite located in the project tag list staged in the local snapshot; invoked at proof line 4830 of the counted proof. It is retained only as a navigation pointer; every load-bearing step is inside the delivered span.\par\medskip
\ifdefined\hypertarget\hypertarget{sourcepointer-7}{}\fi
\noindent\textit{Source reference.} Original lemma 08E5: sources/stacks-chow/tags.txt, source lines 10661--10661. Cross-chapter standard prerequisite located in the project tag list staged in the local snapshot; invoked at proof line 4831 of the counted proof. It is retained only as a navigation pointer; every load-bearing step is inside the delivered span.\par\medskip
\ifdefined\hypertarget\hypertarget{sourcepointer-8}{}\fi
\noindent\textit{Source reference.} Original lemma 08E8: sources/stacks-chow/tags.txt, source lines 10664--10664. Cross-chapter standard prerequisite located in the project tag list staged in the local snapshot; invoked at proof line 4832 of the counted proof. It is retained only as a navigation pointer; every load-bearing step is inside the delivered span.\par\medskip
\ifdefined\hypertarget\hypertarget{sourcepointer-9}{}\fi
\noindent\textit{Source reference.} Original lemma 0A6H: sources/stacks-chow/tags.txt, source lines 12845--12845. Cross-chapter standard prerequisite located in the project tag list staged in the local snapshot; invoked at proof line 4833 of the counted proof. It is retained only as a navigation pointer; every load-bearing step is inside the delivered span.\par\medskip
\ifdefined\hypertarget\hypertarget{sourcepointer-10}{}\fi
\noindent\textit{Source reference.} Original lemma 0AU1: sources/stacks-duality/pinned-original-source.tex, source lines 4686--4686. Same-chapter pointer into the author environment lemma at source line 4686; invoked at proof line 4840 of the counted proof. It is retained only as a navigation pointer; every load-bearing step is inside the delivered span.\par\medskip
\ifdefined\hypertarget\hypertarget{sourcepointer-11}{}\fi
\noindent\textit{Source reference.} Original lemma 0AU0: sources/stacks-duality/pinned-original-source.tex, source lines 4568--4568. Same-chapter pointer into the author environment lemma at source line 4568; invoked at proof line 4843 of the counted proof. It is retained only as a navigation pointer; every load-bearing step is inside the delivered span.\par\medskip
\ifdefined\hypertarget\hypertarget{sourcepointer-12}{}\fi
\noindent\textit{Source reference.} Original lemma 0A9U: sources/stacks-duality/pinned-original-source.tex, source lines 3306--3306. Same-chapter pointer into the author environment lemma at source line 3306; invoked at proof line 4845 of the counted proof. It is retained only as a navigation pointer; every load-bearing step is inside the delivered span.\par\medskip
\ifdefined\hypertarget\hypertarget{sourcepointer-13}{}\fi
\noindent\textit{Source reference.} Original lemma 04II: sources/stacks-chow/tags.txt, source lines 5904--5904. Cross-chapter standard prerequisite located in the project tag list staged in the local snapshot; invoked at proof line 4863 of the counted proof. It is retained only as a navigation pointer; every load-bearing step is inside the delivered span.\par\medskip
\ifdefined\hypertarget\hypertarget{sourcepointer-14}{}\fi
\noindent\textit{Source reference.} Original lemma 068D: sources/stacks-chow/tags.txt, source lines 8002--8002. Cross-chapter standard prerequisite located in the project tag list staged in the local snapshot; invoked at proof line 4867 of the counted proof. It is retained only as a navigation pointer; every load-bearing step is inside the delivered span.\par\medskip
\ifdefined\hypertarget\hypertarget{sourcepointer-15}{}\fi
\noindent\textit{Source reference.} Original lemma 0AA2: sources/stacks-duality/pinned-original-source.tex, source lines 4624--4624. Same-chapter pointer into the author environment lemma at source line 4624; invoked at proof line 4875 of the counted proof. It is retained only as a navigation pointer; every load-bearing step is inside the delivered span.\par\medskip
\ifdefined\hypertarget\hypertarget{sourcepointer-16}{}\fi
\noindent\textit{Source reference.} Original lemma 0AA1: sources/stacks-duality/pinned-original-source.tex, source lines 4609--4609. Same-chapter pointer into the author environment lemma at source line 4609; invoked at proof line 4883 of the counted proof. It is retained only as a navigation pointer; every load-bearing step is inside the delivered span.\par\medskip
\begin{lemma}
\label{lemma-perfect-comparison-shriek}
In Situation \ref{situation-shriek} let $f : X \to Y$ be a morphism of
$\textit{FTS}_S$. Assume $f$ is perfect (e.g., flat). Then
\begin{enumerate}
\item[(a)] $f^!\mathcal{O}_Y$ is in $D_{\textit{Coh}}^b(\mathcal{O}_X)$,
\item[(b)] $f^!\mathcal{O}_Y$ has finite tor dimension in
$D(f^{-1}\mathcal{O}_Y)$,
\item[(c)] $\mathcal{O}_X \to
R\SheafHom_{\mathcal{O}_X}(f^!\mathcal{O}_Y, f^!\mathcal{O}_Y)$
is an isomorphism,
\item[(d)] $f^!$ maps $D_{\textit{Coh}}^b(\mathcal{O}_Y)$ into
$D_{\textit{Coh}}^b(\mathcal{O}_X)$,
\item[(e)] the map
$\mu_{f,  K} :
Lf^*K \otimes_{\mathcal{O}_X}^\mathbf{L} f^!\mathcal{O}_Y
\to
f^!K$
of Lemma \ref{lemma-map-pullback-to-shriek-well-defined}
is an isomorphism for all $K \in D_\QCoh^+(\mathcal{O}_Y)$.
\end{enumerate}
\end{lemma}

\begin{proof}
(A flat morphism of finite presentation is perfect, see
More on Morphisms, Lemma
\ref{more-morphisms-lemma-flat-finite-presentation-perfect}.)
Assertions (a), (b), and (c) are local on $X$.
Thus we may assume $X$ and $Y$ are affine.
Then Remark \ref{remark-local-calculation-shriek}
turns (a), (b), and (c) into (1)(a), (1)(b), and (1)(c) of
Dualizing Complexes, Lemma \ref{dualizing-lemma-relative-dualizing-algebraic}.
(Use Derived Categories of Schemes, Lemmas
\ref{perfect-lemma-tor-dimension-rel-affine},
\ref{perfect-lemma-pseudo-coherent},
\ref{perfect-lemma-identify-pseudo-coherent-noetherian}
\ref{perfect-lemma-quasi-coherence-internal-hom} to match the assertions.)

\medskip\noindent
To prove (d) and (e) we begin with a series of preliminary remarks.
\begin{enumerate}
\item We already know that $f^!$ sends $D_{\textit{Coh}}^+(\mathcal{O}_Y)$
into $D_{\textit{Coh}}^+(\mathcal{O}_X)$, see
Lemma \ref{lemma-shriek-coherent}.
\item If $f$ is an open immersion, then $f^! = f^*$ and (d) and (e)
are true because we can take $\overline{X} = Y$ in the construction
of $f^!$ and $\mu_f$, see Lemma \ref{lemma-shriek-open-immersion}.
\item If $f$ is a perfect proper morphism, then (e) is true by
Lemma \ref{lemma-compare-with-pullback-flat-proper-noetherian}.
\item If there exists an open covering $X = \bigcup U_i$ and (d) is
true for $U_i \to Y$, then (d) is true for $X \to Y$. Same for (e).
This holds because the construction of $f^!$ and $\mu_f$ commutes
with passing to open subschemes.
\item If $g : Y \to Z$ is a second perfect morphism in $\textit{FTS}_S$
and (e) holds for $f$ and $g$, then
$f^!g^!\mathcal{O}_Z =
Lf^*g^!\mathcal{O}_Z \otimes_{\mathcal{O}_X}^\mathbf{L} f^!\mathcal{O}_Y$
and (e) holds for $g \circ f$ by the commutative diagram
of Lemma \ref{lemma-map-pullback-to-shriek-well-defined}.
\item If (d) and (e) hold for both $f$ and $g$, then
(d) and (e) hold for $g \circ f$. Namely, then $f^!g^!\mathcal{O}_Z$
is bounded above (by the previous point) and $L(g \circ f)^*$ has finite
cohomological dimension and (d) follows from (e) which we saw above.
\end{enumerate}
From these points we see it suffices to prove (d) and (e) in case $X$ is
affine. Choose an immersion $X \to \mathbf{A}^n_Y$
(Morphisms, Lemma \ref{morphisms-lemma-quasi-affine-finite-type-over-S})
which we factor as $X \to U \to \mathbf{A}^n_Y \to Y$ where $X \to U$
is a closed immersion and $U \subset \mathbf{A}^n_Y$ is open.
Note that $X \to U$ is a perfect closed immersion by
More on Morphisms, Lemma \ref{more-morphisms-lemma-perfect-permanence}.
Thus it suffices to prove the lemma for a perfect closed immersion
and for the projection $\mathbf{A}^n_Y \to Y$.

\medskip\noindent
Let $f : X \to Y$ be a perfect closed immersion. We already know (d) holds.
Let $K \in D^b_{\textit{Coh}}(\mathcal{O}_Y)$.
Then $f^!K = R\SheafHom(\mathcal{O}_X, K)$
(Lemma \ref{lemma-shriek-closed-immersion})
and $f_*f^!K = R\SheafHom_{\mathcal{O}_Y}(f_*\mathcal{O}_X, K)$.
Since $f$ is perfect, the complex $f_*\mathcal{O}_X$ is perfect
and hence $R\SheafHom_{\mathcal{O}_Y}(f_*\mathcal{O}_X, K)$ is bounded above.
This proves that (d) holds. Some details omitted.

\medskip\noindent
Let $f : \mathbf{A}^n_Y \to Y$ be the projection. Then (d) holds
by repeated application of Lemma \ref{lemma-shriek-affine-line}.
Finally, (e) is true because it holds for $\mathbf{P}^n_Y \to Y$
(flat and proper) and because $\mathbf{A}^n_Y \subset \mathbf{P}^n_Y$
is an open.
\end{proof}

\end{document}
